\documentclass[12pt,a4paper]{report}

\usepackage[utf8]{inputenc}
\usepackage{graphicx}
\usepackage{svg}
\usepackage{geometry}
\geometry{a4paper, margin=1in}
\usepackage{titlesec}
\usepackage{hyperref}
\usepackage{booktabs}
\usepackage{longtable}
\usepackage{float}
\usepackage{listings}
\usepackage{xcolor}

\titleformat{\chapter}[display]
  {\normalfont\huge\bfseries}{\chaptertitlename\ \thechapter}{20pt}{\Huge}

\begin{document}

\begin{titlepage}
    \centering
    \vspace*{1cm}
    
    {\Large \textsc{B.Sc. Engg. Project}}\\[1.5cm]
    
    {\huge \textbf{SELLWISE (AI Sales Analytics)}}\\[1.5cm]
    
    {\Large Developed by}\\[0.5cm]
    {\large Student Name (ID: 123456789)}\\[1.5cm]
    
    {\Large Supervised by}\\[0.5cm]
    {\large Supervisor Name}\\[0.2cm]
    {\large Lecturer}\\[1.5cm]
    
    {\Large Department of Computer Science \& Engineering}\\[0.5cm]
    {\large (In partial fulfillment of the requirements for the degree of\\
    Bachelor of Science in Computer Science \& Engineering)}\\[1.5cm]
    
    \vfill
    
    {\Large Department of Computer Science \& Engineering}\\[0.5cm]
    {\Large Bangladesh University of Business \& Technology (BUBT)}\\[0.5cm]
    {\Large Dhaka 1000}\\[0.5cm]
    {\large \today}
\end{titlepage}

\pagenumbering{roman}
\tableofcontents
\newpage
\listoffigures
\newpage
\pagenumbering{arabic}

\chapter{INTRODUCTION}
\section{Introduction}
SellWise is a web-based application built to simplify the daily demands of running a retail business. It brings sales operations, stock management, and customer tracking together into one unified platform. Most small shops today still rely on paper notebooks or basic spreadsheets to record their sales and inventory — a manual process that is slow, error-prone, and hard to scale. SellWise acts as a digital backbone that keeps everything organized and within reach.

Beyond day-to-day operations, the system uses machine learning to forecast product demand for the coming weeks. The project is split into two components: a main web application built with ASP.NET Core that handles all transactions and user interactions, and a separate Python service responsible for running the sales predictions. Separating these two concerns keeps the web application lean and responsive, regardless of forecast computation time.

\section{Project overview}
This project provides an affordable and easy-to-use software solution for small and medium-sized shops. Users don't need any formal technical knowledge to use this system. The website is designed to be error-free, secure, and reliable. 

Each shop has its own isolated workspace, meaning data from one store never mixes with another. The web application handles the transactions, and the background machine learning service handles the predictions. This way, shop owners can focus on improving their business rather than getting stuck with record keeping.

\section{Problem Statement}
Currently, many small businesses face several challenges because they don't have proper management software. The current situation is very limited and often leads to the following problems:
\begin{itemize}
    \item \textbf{Inventory Issues:} Shops often overstock items that don't sell, or they run out of popular items and miss out on sales.
    \item \textbf{Losing Customers:} Since shop owners don't track customer purchase history properly, they don't realize when a regular customer stops coming.
    \item \textbf{Scattered Information:} Sales are noted in one book, inventory in another, making it hard to get a clear picture of profit or loss.
    \item \textbf{No Future Planning:} Owners can only see what happened in the past, but they don't have tools to predict what they will need next month.
\end{itemize}

\section{Motivation}
Facing the problems described above, we wanted to build a custom system called SellWise that can solve these challenges. Our motivation comes from the desire to:
\begin{itemize}
    \item Give shop owners a clear view of their business from a single dashboard.
    \item Help them prepare properly by forecasting future product demand.
    \item Track customer buying habits to offer better services.
    \item Provide an inexpensive, efficient, and comfortable way to manage a shop.
    \item Prevent data loss and calculation mistakes that happen with manual paperwork.
\end{itemize}

\section{Objectives}
\begin{itemize}
    \item Provide a user-friendly environment for managing shop inventory and sales.
    \item Implement a machine learning model to predict product demand for the next 30 days.
    \item Ensure all sales transactions are processed safely without corrupting inventory data.
    \item Group customers automatically based on how frequently they buy.
    \item Provide an alert system that warns the owner when product stock is low.
    \item Generate downloadable PDF reports for daily or monthly sales.
    \item Reduce paperwork and manual calculations.
\end{itemize}

\section{Development Tools}
SellWise is a web application integrated with a relational database and a standalone machine learning service. The core languages used across both components are C\#, Python, HTML, CSS, and JavaScript.
\begin{itemize}
    \item \textbf{Technology:}
    \begin{itemize}
        \item SQL Server for database management.
        \item ASP.NET Core MVC for the main web framework.
        \item FastAPI (Python) for the machine learning service.
        \item Meta Prophet for time-series forecasting.
        \item Docker for containerizing the database.
    \end{itemize}
    \item \textbf{Language:}
    \begin{itemize}
        \item C\#
        \item Python
        \item HTML
        \item CSS
        \item JavaScript
        \item SQL
    \end{itemize}
\end{itemize}

\section{Proposed System}
In our proposed system, there are three types of actors: Business Owner, Store Manager, and Employee.
Users will log in or register themselves using their email and password. A single admin account can manage multiple stores. 

The Owner has full control over the system. The Manager can process orders, check inventory, and view analytics. The Employee can only process orders and check the product catalog. 
During a sale, the system automatically checks if there is enough stock. If a product is running low, the system will show an alert. The system will also study the past sales data to forecast future demand, which will be shown on the owner's dashboard.

\section{Purpose Of The System}
Many small shops know they should keep better track of their sales, but they can't afford expensive software. The purpose of SellWise is to encourage business owners to adopt digital tools, no matter where they are. 

By bringing everything into one system, shop owners don't have to jump between different tools. They can see all their orders, inventory, and reports in one place. The system automatically does the hard work, like calculating which customers are loyal and predicting future sales. This will provide opportunities for small businesses to grow and achieve their goals.

\section{Goal}
The major goal of this project is to provide a user-friendly environment to manage retail shops efficiently. Another goal is to enhance the quality of business decisions by providing accurate data and predictions. We also aim to improve user accessibility and time flexibility so that store owners can check their business status from anywhere.

\chapter{EXISTING SYSTEM}

\section{Introduction}
Technology is among the most important factors for business growth today. Although many shops still use manual methods, the use of software increases efficiency and reduces errors. There are two main aims to this existing system analysis. The first aim is to establish the characteristics and importance of retail software. The second aim is to identify the techniques, tools, and approaches used in existing systems, to discuss their effectiveness, and to uncover factors that limit their usage by small shops. 

\section{Existing System}
We have visited and analyzed several existing systems. They are listed below.

\subsection{Shopify Analytics}
Shopify is one of the most popular e-commerce platforms that provides built-in analytics features. It offers basic sales reporting, traffic analytics, and inventory management.

\textbf{Advantages:}
\begin{itemize}
    \item User-friendly interface with intuitive dashboards.
    \item Integrated e-commerce and POS capabilities.
    \item Basic inventory tracking and low-stock alerts.
\end{itemize}

\textbf{Disadvantages:}
\begin{itemize}
    \item Limited advanced analytics without expensive add-ons.
    \item No built-in machine learning or predictive analytics.
    \item Advanced features require third-party app integrations at additional cost.
\end{itemize}

\subsection{Square for Retail}
Square provides a free POS system with built-in inventory management and basic analytics for small retail businesses.

\textbf{Advantages:}
\begin{itemize}
    \item Free base POS system with inventory tracking.
    \item Basic sales analytics and reporting.
    \item Customer directory with purchase history.
\end{itemize}

\textbf{Disadvantages:}
\begin{itemize}
    \item Advanced analytics require a paid Square Analytics subscription.
    \item No predictive demand forecasting.
    \item Restrictive for multi-store operations.
\end{itemize}

\subsection{Lightspeed Retail}
Lightspeed offers a comprehensive retail POS system with advanced inventory management and analytics features.

\textbf{Advantages:}
\begin{itemize}
    \item Advanced inventory management with multi-location support.
    \item Detailed sales analytics and reporting.
    \item Vendor management and purchase ordering.
\end{itemize}

\textbf{Disadvantages:}
\begin{itemize}
    \item Higher price point compared to competitors.
    \item No built-in machine learning for demand prediction.
    \item Complex setup and configuration required.
\end{itemize}

\section{Supporting Literature}

\subsection{Used Diagram}

\textbf{Entity Relationship Diagram}\\
An Entity-relationship model (ER model) describes the structure of a database with the help of a diagram, which is known as Entity Relationship Diagram (ER Diagram). An ER model is a design or blueprint of a database that can later be implemented as a database. The main components of E-R model are:
\begin{itemize}
    \item \textbf{Entity:} An entity set is a group of similar entities and these entities can have attributes. In terms of DBMS, an entity is a table.
    \item \textbf{Attribute:} Attributes are the properties which define the entity type.
    \item \textbf{Relationship:} A relationship is represented by a diamond shape in the ER diagram; it shows the relationship among entities.
\end{itemize}

\textbf{Data Flow Diagram}\\
A data-flow diagram is a way of representing a flow of data through a process or a system. The DFD also provides information about the outputs and inputs of each entity and the process itself. Components of DFD involve:
\begin{itemize}
    \item \textbf{Process:} An activity that changes or transforms data flows.
    \item \textbf{Data Flow:} Movement of data between external entities, processes, and data stores.
    \item \textbf{Data Store:} A data store simply holds data for later access.
    \item \textbf{External entity:} Also known as actors, sources, or sinks.
\end{itemize}

\textbf{Use Case Diagram}\\
A use case diagram is a way to summarize details of a system and the users within that system. It is generally shown as a graphic depiction of interactions among different elements in a system. 
\begin{itemize}
    \item \textbf{System:} A specific sequence of actions and interactions between actors and the system.
    \item \textbf{Actors:} The users that interact with a system.
\end{itemize}

\subsection{Technology Used}
To build this system, we used ASP.NET Core, HTML, CSS, and JavaScript for the web layer, with SQL Server as the underlying database.
\begin{itemize}
    \item \textbf{ASP.NET Core MVC:} Used for the web application framework. It keeps data (Model), user interface (View), and request handling (Controller) separate.
    \item \textbf{Entity Framework Core:} Used to talk to the SQL Server database. Instead of writing raw SQL, we write queries in C\#.
    \item \textbf{Python FastAPI:} Used to build the machine learning service for sales prediction.
    \item \textbf{Facebook Prophet:} An open-source time-series forecasting library used to predict product demand.
    \item \textbf{Bootstrap \& Chart.js:} Used for designing the frontend and drawing graphs on the dashboard.
\end{itemize}

\section{Conclusion}
Software systems are becoming increasingly prominent in retail management. However, we found that most retail analytics tools are either too basic or too heavy. SellWise tries to solve this by providing built-in forecasting, automatic customer grouping, and a multi-store architecture in a simple, easy-to-use package.


\chapter{SYSTEM EVALUATION}

\section{Introduction}
Research in retail management covers a wide range of operations, from point-of-sale systems to remote reporting and analytics. Given the clear gaps in existing basic software, we designed SellWise using modern web frameworks and machine learning — combining sales tracking, real-time inventory alerts, and demand forecasting into a single, coherent system. The goal is to give store owners tools that are both practical to use and genuinely helpful for making better business decisions.

\section{Requirement Analysis}
Requirement Analysis is the process of defining user expectations for a new software being built. It encompasses those tasks that go into determining the needs or conditions to meet for a new product.

\subsection{Non-Functional requirement}
Types of Non-functional Demand:
\begin{itemize}
    \item \textbf{Performance:} Dashboard should load quickly, and operations should complete within seconds.
    \item \textbf{Security Demand:} Passwords must be hashed, and users can only access features based on their roles.
    \item \textbf{Reliability:} Data must not be corrupted during a transaction failure.
    \item \textbf{Manageability demand:} Easy to maintain and update the code.
    \item \textbf{Usability demand:} The website must be responsive and mobile-friendly.
\end{itemize}

\subsection{Functional requirement}
A Functional Requirement describes the service that the software must offer.
\begin{itemize}
    \item User Sign Up and Sign In.
    \item Create and manage products.
    \item Add products to the cart and process orders.
    \item View sales dashboard and statistics.
    \item View predicted demand for products.
    \item View low-stock alerts.
    \item Manage customer profiles.
    \item Generate PDF sales reports.
\end{itemize}

\section{Methodology}

\subsection{Agile Development Model}
We adopted the Agile Development Model for this project. Agile is an iterative, people-focused methodology that emphasizes delivering working software in short cycles and continuously improving based on feedback.

\textbf{Reasons for choosing Agile Development Model:}
\begin{itemize}
    \item Features are delivered incrementally, which allows for early feedback and course correction.
    \item The development process remains flexible enough to accommodate changing requirements.
    \item Team communication is more frequent and transparent.
    \item Defects are identified and resolved early in the development cycle, rather than at the end.
\end{itemize}

\subsection{Gantt Chart}

\begin{figure}[H]
    \centering
    \includesvg[width=\textwidth]{"diagrams/Gantt Chart Diagram.drawio.svg"}
    \caption{Gantt Chart}
\end{figure}

\begin{table}[H]
\centering
\begin{tabular}{|l|c|c|c|c|c|}
\hline
\textbf{Task} & \textbf{Week 1-2} & \textbf{Week 3-4} & \textbf{Week 5-6} & \textbf{Week 7-8} & \textbf{Week 9-10} \\
\hline
Database \& Auth & xxxx & & & & \\
\hline
Products \& POS & & xxxx & & & \\
\hline
Dashboard \& ML & & & xxxx & & \\
\hline
RFM \& Alerts & & & & xxxx & \\
\hline
Testing \& Polish & & & & & xxxx \\
\hline
\end{tabular}
\caption{Project Timeline}
\end{table}

\subsection{Data Flow Diagram}

\begin{figure}[H]
    \centering
    \includesvg[width=\textwidth]{"diagrams/DFD LEVEL0 Diagram.drawio.svg"}
    \caption{Context Level DFD (Level 0)}
\end{figure}

\begin{figure}[H]
    \centering
    \includesvg[width=\textwidth]{"diagrams/DFD LEVEL 1 Diagram.drawio.svg"}
    \caption{Level 1 Data Flow Diagram}
\end{figure}

\subsection{ER Diagram}

\begin{figure}[H]
    \centering
    \includesvg[width=\textwidth]{"diagrams/ER Diagram.drawio.svg"}
    \caption{ER Diagram}
\end{figure}

\textbf{Key Entities:}
\begin{itemize}
    \item User, Store, StoreMember
    \item Product, Order, OrderItem
    \item Customer, Expense, InventoryAlert, Forecast
\end{itemize}

\subsection{Use Case Diagram}

\begin{figure}[H]
    \centering
    \includesvg[width=\textwidth]{"diagrams/USE CASE Diagram.drawio.svg"}
    \caption{Use Case Diagram}
\end{figure}

\begin{itemize}
    \item \textbf{Owner:} Manage Staff, View Dashboard, Switch Store, Generate Reports.
    \item \textbf{Manager:} Process Orders, Manage Inventory, View Analytics, Manage Customers.
    \item \textbf{Employee:} Process Orders, View Catalog.
\end{itemize}

\subsection{Activity Diagram}

\begin{figure}[H]
    \centering
    \includesvg[width=\textwidth]{"diagrams/New_06_Activity_Diagram.svg"}
    \caption{Activity Diagram}
\end{figure}

Order Creation Activity:\\
User clicks "Create Order" $\rightarrow$ Adds items to cart $\rightarrow$ Selects customer $\rightarrow$ Submits form $\rightarrow$ System begins database transaction $\rightarrow$ Validates stock $\rightarrow$ If sufficient, deducts stock and saves order $\rightarrow$ Commits transaction.

\subsection{Sequence Diagram}

\begin{figure}[H]
    \centering
    \includesvg[width=\textwidth]{"diagrams/new sequence Diagram.drawio.svg"}
    \caption{Sequence Diagram}
\end{figure}

Dashboard Sequence:\\
Browser requests Dashboard $\rightarrow$ AnalyticsService checks Database $\rightarrow$ Checks Forecast Cache $\rightarrow$ If miss, requests Python ML Service $\rightarrow$ Python returns prediction $\rightarrow$ AnalyticsService saves cache $\rightarrow$ Returns data to Browser.

\subsection{Class Diagram}

\begin{figure}[H]
    \centering
    \includesvg[width=\textwidth]{"diagrams/Class Diagram.drawio.svg"}
    \caption{Class Diagram}
\end{figure}

\begin{itemize}
    \item \textbf{BaseController:} Handles current store session.
    \item \textbf{DashboardController, OrderController:} Inherit BaseController.
    \item \textbf{AnalyticsService:} Calculates KPIs.
    \item \textbf{ForecastService:} Talks to the Python ML API.
\end{itemize}

\subsection{Database Design}
Important tables:
\begin{itemize}
    \item \textbf{Products:} Stores product details and stock quantity.
    \item \textbf{Orders:} Stores transaction totals and dates.
    \item \textbf{OrderItems:} Links orders with products and quantities.
    \item \textbf{Customers:} Stores customer details and RFM scores.
\end{itemize}


\chapter{DESIGN IMPLEMENTATION}

\section{User Section}
We have three kinds of users:
\begin{itemize}
    \item Owner
    \item Manager
    \item Employee
\end{itemize}

We will show the background implementation of various features in this part.

\subsection{Database Connection}
The SQL Server connection is configured in \texttt{appsettings.json} and used by Entity Framework Core to connect to the containerized SQL Server instance.
\begin{lstlisting}[language=json, basicstyle=\small\ttfamily, backgroundcolor=\color{lightgray!20}]
{
  "ConnectionStrings": {
    "DefaultConnection": "Server=localhost,1433;Database=SellWise;
    User Id=sa;Password=YourPass123!;TrustServerCertificate=True;"
  }
}
\end{lstlisting}

\subsection{Authentication}
The login action authenticates the user via ASP.NET Core Identity, then sets the active store ID in the session to scope all subsequent queries to that store.

\subsection{POS Order Processing}
When a user creates an order, the system opens a database transaction to ensure that stock validation and deduction either both succeed or both roll back atomically. This prevents any scenario where payment is recorded but stock is not deducted, or vice versa. The lockout policy also protects the system from brute-force login attempts.

\subsection{Inventory Alerts}
The system automatically scans all active products and generates low-stock alerts when a product's quantity falls at or below its configured threshold. Products that recover above the threshold are automatically resolved.

\subsection{Customer Segmentation (RFM)}
Customers are automatically scored and segmented using the RFM model — Recency (how recently they bought), Frequency (how often), and Monetary value (how much they spent). 

\subsection{Machine Learning Forecasting}
The Python microservice uses Meta's Prophet library to generate a 30-day demand forecast per product. The model enables weekly seasonality by default, and only enables yearly seasonality when the historical data spans at least one full year — avoiding overfitting on short datasets.

\section{Conclusions}
Several design decisions made during implementation proved particularly effective:
\begin{enumerate}
    \item Encapsulating business logic within service classes kept controllers thin and the codebase easy to maintain.
    \item Using database transactions for order processing eliminated any risk of inventory mismatches during concurrent operations.
    \item Filtering all queries by Store ID at the service layer ensured strict data isolation between different tenants.
    \item Implementing a fallback to moving average ensures the dashboard remains fully functional even when the ML service is temporarily unreachable.
\end{enumerate}


\chapter{TESTING}

\section{Testing}
Testing is the process in which the system is run on manually created input so that we can verify the system is correctly working as desired or not. During systems testing, the system is used experimentally to ensure that the software does not fail. Special test data are input for processing, and the results examined.

\section{Validation Criteria}
The validation criteria in this project are as follows:
\begin{itemize}
    \item Here, the user inputs are validated before storing them.
    \item All the screens have a similar look and feel, providing a better user interface.
    \item Transactions cannot corrupt data if an error occurs mid-process.
\end{itemize}

\section{Importance of Testing}
The importance of system testing is that the system is expected to run according to requirement before delivering it to the user. The System is tested on the basis of specification so that it does not fail on the user site.

\section{Structural Testing}
Structural Testing takes into account the internal mechanism of a system or component. We checked the code coverage and error handling, ensuring that all database operations are wrapped in try-catch blocks.

\section{Functional Testing}
We tested each feature against what it is supposed to do, such as login, adding items to the cart, processing orders, and checking whether the dashboard numbers match the database.

\section{Acceptance Test Generation}
The objective of this step is to produce a set of test data that may be used to test the system. We generated demo data with 2 stores, hundreds of customers, and 180 days of order history to simulate a real environment.

\section{Unit Testing}
Unit testing focuses on the modules independently to locate errors. We checked individual service methods like the RFM calculation and the Prophet forecasting logic to ensure they return accurate numbers.

\section{Integration Testing}
This is a systematic technique for constructing the program structure while uncovering errors associated with interfaces. We verified the connection between the C\# web application and the Python machine learning service.

\section{Test Data Used}
\begin{itemize}
    \item \textbf{Using Live Test Data:} Live tests were conducted by inputting real-world retail scenarios.
    \item \textbf{Using Artificial Test Data:} Demo scripts were used to inject bulk data into the database for load testing.
\end{itemize}

\section{Test Cases}
\begin{table}[H]
\centering
\begin{tabular}{|l|p{6cm}|l|l|}
\hline
\textbf{Module} & \textbf{Test Case Scenarios} & \textbf{Tested} & \textbf{Working} \\
\hline
Authentication & 1. Verify login with valid credentials\newline 2. Verify account lockout\newline 3. Verify signup creates store & Yes & Yes \\
\hline
Dashboard & 1. Verify KPI calculations\newline 2. Verify charts render correctly & Yes & Yes \\
\hline
POS Orders & 1. Verify cart adds items\newline 2. Verify stock validation\newline 3. Verify transaction rollback on failure & Yes & Yes \\
\hline
Products & 1. Verify CRUD operations\newline 2. Verify CSV import & Yes & Yes \\
\hline
Customers & 1. Verify RFM calculation\newline 2. Verify segment assignment & Yes & Yes \\
\hline
\end{tabular}
\caption{Test Cases}
\end{table}

\section{Testing result}
Sign in system works perfectly. If the given password doesn't match with the database, then it shows a warning. Order creation works smoothly and safely deducts inventory.

\section{Testing Summary}
We tested all the panels assigned for users. The project was tested in the localhost with all possible options. The testing process involved successful login, form validation, order creation, and record generation. The project was successfully tested and a few bugs were fixed.


\chapter{USER MANUAL}

\section{Introduction}
The system provides the capability to manage all aspects of your retail business by providing a complete management suite. By creating powerful resources and delivering them online, the system is flexible and easy to use.

\section{System Requirements}

\subsection{H/W Requirements}
\begin{itemize}
    \item \textbf{Operating System:} Windows 10, macOS, Linux
    \item \textbf{Processor:} Dual core processor
    \item \textbf{RAM:} Minimum 4GB
    \item \textbf{Storage:} 128GB
    \item \textbf{Network:} Internet connection required.
\end{itemize}

\subsection{S/W Requirements}
\begin{itemize}
    \item \textbf{Browser:} Chrome, Firefox, or Edge.
    \item \textbf{Docker:} Docker Desktop for running the database.
    \item \textbf{Python:} Python 3.11+ for the ML service.
\end{itemize}

\section{User Interface}
The very first portion is "User Login" which is the user interface to access the system. Users must validate their email and password.
After login, users see the Dashboard where they can view sales KPIs, graphs, and demand forecasts. 
To process a sale, the user goes to the "Orders" section, adds products to the cart, applies any discounts, and submits the order.
To manage inventory, the user goes to the "Products" page to add or edit items.
The "Customers" page shows all customer profiles along with their automatically generated loyalty segments.

\section{Conclusion}
In this chapter, we described the requirement specifications of the system and the actions that can be done. User will discover the basic characteristics and the way to use them easily.


\chapter{CONCLUSION AND FUTURE WORKS}

\section{Conclusion}
SellWise was built to address a genuine gap that many small retail businesses face — the lack of affordable, integrated tools for managing sales, inventory, and customers. Over the course of this project, we successfully developed a full-stack application that combines a point-of-sale system, real-time inventory management, machine learning-based demand forecasting, and automatic customer segmentation into a single, cohesive platform.

The experience reinforced how important it is to design with real users in mind. By keeping the interface simple and the logic well-separated, we were able to deliver a system that is both technically sound and straightforward to use in practice.

\section{Outcomes}
\begin{itemize}
    \item Shop owners can now manage their entire business from a single dashboard, saving significant time compared to manual methods.
    \item Data-driven demand forecasting helps avoid both overstocking and stockouts.
    \item Automatic RFM-based customer segmentation helps identify loyal customers and those at risk of leaving.
    \item Atomic order processing ensures daily sales are recorded accurately without data corruption.
    \item Low-stock alerts notify the owner before a product runs out, reducing lost sales.
\end{itemize}

\section{Current Limitations}
The system works well in its current state, but there are a few constraints worth noting:
\begin{itemize}
    \item The ML forecasting model is limited to Prophet; it does not automatically compare or switch between alternative models.
    \item The system requires an active internet or local network connection — there is currently no offline mode.
    \item Each store supports only a single currency, which limits use in multi-currency markets.
\end{itemize}

\section{Future Works}
Several improvements are planned to extend the platform's capabilities:
\begin{itemize}
    \item \textbf{Email Marketing Integration:} Connect customer segments directly to email marketing platforms to enable automated, targeted campaigns.
    \item \textbf{Mobile Application:} Develop a native mobile app to allow store owners and managers to run operations on the go, with offline support for low-connectivity environments.
    \item \textbf{Cloud Deployment:} Migrate the infrastructure to a cloud provider such as AWS or Azure to improve availability, scalability, and ease of deployment for end users.
\end{itemize}

\end{document}
